\documentclass[12pt]{article}
\usepackage[utf8]{inputenc}
\usepackage[T1]{fontenc}
\usepackage{hyperref}
\title{San Luís Temalacayuca Popoloca Language Resources for K-12 Curriculum Development}
\author{Compiled by Claude Sonnet 5.5 (claude-sonnet-5-5) for Andr\'as Kornai}
\date{}
\begin{document}
\maketitle
\section{Introduction}
San Luís Temalacayuca Popoloca (ISO 639-3 pps) is an Oto-Manguean (Popolocan) language of the Tehuacán valley in southern Puebla, Mexico. Speakers call it Ngigua. It is closely related to San Marcos Tlacoyalco Popoloca (pls); sources report over 90\% mutual intelligibility. Speaker figures in the sources found are about 4,700 for this variety. Search results did not reveal any school curriculum in the language.
\section{Primary Educational Resources}
No primers, textbooks or school programmes for pps were found in the search results. Any curriculum will have to be built from the reference and scripture materials below.
\section{Linguistic and Reference Materials}
\begin{itemize}
\item Wikipedia, Northern Popoloca language (covers pps and pls together): \url{https://en.wikipedia.org/wiki/Northern_Popoloca_language}
\item Omniglot, Northern Popoloca (Ngigua) overview: \url{https://www.omniglot.com/writing/northernpopoloca.htm}
\item Tlalocan journal article (UNAM) that appears to contain text and grammatical notes on the variety; content not verified beyond the search snippet: \url{https://revistas-filologicas.unam.mx/tlalocan/index.php/tl/article/download/444/435/661}
\item ASJP language entry with a word list: \url{https://asjp.clld.org/languages/SAN_LUIS_TEMALACAYUCA_POPOLOCA}
\item Wikitongues language page: \url{https://wikitongues.org/languages/pps/}
\end{itemize}
\section{Religious and Community Resources}
\begin{itemize}
\item Temalacayuca Popoloca New Testament (Digital Bible Society): \url{https://dbs.org/bibles/PPSWBT}
\end{itemize}
\section{Digital and Audio Resources}
\begin{itemize}
\item Global Recordings Network, Popoloca, San Luis Temalacayuca (audio recordings): \url{https://www.globalrecordings.net/en/language/pps}
\end{itemize}
\section{Curriculum Development Context}
The language appears vigorous, and an alphabet is reported to have been established in 2005. No evidence of school use was found. Community contact, orthography confirmation and Mexican intercultural bilingual education bodies would be the next step; none are cited here because none were found in the searches.
\end{document}